% Transcription — fonds Alexandre Grothendieck, shelfmark 132, batch 1 (pages 1-20).
% Established from the facsimile archives/batches/132/batch-01.pdf, cut from a
% copy of 132.pdf fetched through the project's own relay
% (https://grothendieck-relay.onrender.com/source/132.pdf), not uploaded by the
% user, per the skill transcribe-grothendieck, with the 700 dpi tiles of
% archives/tiles/132/ (margins turned upright for pages 10 and 19). The batch
% file carries no cover sheet: PDF page k is page k of the folder, and the
% archivists' pencil numbers bottom-left agree on every leaf. The folder is
% thirty-one pages; this is the first of two batches (pages 21-31 are batch 2,
% transcribed in parallel by another pass and already in the repository; the
% notation below was aligned with it).
%
% Pass: Opus 5.5 (claude-opus-5-5), 2026-09-26 — first pass, unchecked against the pages by a human.
%
% Dating. The inventory dates the folder « 1974 » (copied verbatim into
% \dating, as in batch 2), which is the date of the offprints. Since the folder
% has a date of its own, the group rule does not apply; for the record, the
% inventory group « Descentes » (126-128, which also takes 129-133) is dated
% « [avant 1970] », which does not agree. The manuscript leaves carry no date.
%
% Folder context. Folder 132, « Plongements dans surfaces (Ladegaillerie) :
% notes manuscrites (s.d.), tirés à part (1974) », in the notebook Notes
% techniques dispersées. The offprints (two Notes aux Comptes rendus by Yves
% Ladegaillerie) are all in batch 2 (pages 25-31) and are summarised there;
% no printed leaf falls in this batch, so nothing here is another author's
% text. Page 18 names « Yves » once (« Le calcul des g_i n_i est donné par
% Yves »), which is the only explicit link in this batch to Ladegaillerie.
%
% Pages skipped: 4, 9 and 12, blank sheets (only show-through).
% Pages summarised: none. Page 17 carries, between two horizontal rules, a
% block of scratch diagrams struck through by long diagonals; it is described
% in one \note rather than set, the rest of the page is transcribed.
% Pages transcribed from his hand: 1-3, 5-8, 10, 11, 13-20 (seventeen pages),
% all in blue ink, fast drafting except page 13, which is written almost fair.
%
% His pagination: pages 1-3 are numbered (1), (2), (3) by him; nothing else
% is. His own sectioning letters appear as circled « (A) » (page 16), « (B) »
% (page 19) and « (B') » (page 20); batch 2 continues with (B') on page 21 and
% (C) on page 22. The conditions alpha), beta), gamma) that batch 2 refers to
% are set up on page 19; pages 13-15 have their own « 1) » with sub-cases
% alpha) « Réduction au cas N discret » and beta) « Cas H connexe ».
%
% What the batch is about. Three runs.
%   1-3    Fibre products of spaces and fundamental groupoids: for f : X -> S
%          a Serre fibration and g : Y -> S, the functor
%          Pi_1(X x_S Y) -> Pi_1(X) x^2_{Pi_1(S)} Pi_1(Y) is bijective on
%          pi_0 and surjective on pi_1; the 2-fibre product of Cat G and
%          Cat H over Cat Sigma, its classes = double cosets H\Sigma/G,
%          pi_0(Z) = pi_1(Y)\pi_1(S)/pi_1(X); comparison of the homotopy
%          exact sequences (page 3) giving: Z -> Z~ induces an iso on pi_1
%          iff pi_2(X) x pi_2(Y) -> pi_2(S) is onto, and pi_i(Z) = 0 for
%          i >= 2 when X, Y, S have none.
%   5-11   Groupoid models of surfaces with holes: objects
%          (C, D, phi : D -> C), C a « 2-special » groupoid Cat G, D a sum of
%          copies of Cat Z (boundary loops l_i in G), 1-arrows and 2-arrows;
%          pi_0 and pi_1 of End(C, D, phi): pi_0 is the subgroup of Autext G
%          of classes sending each l_i to a conjugate of some l_j, pi_1 = 0
%          except for the sphere with one or two holes; with D fixed, a
%          « special Teichmüller » group TS and an exact sequence
%          0 -> Z^I -> TS -> T -> S_I -> 1 (page 8), Ker Psi = Z^I except for
%          the cylinder (Z) and the disc. Page 10 sets up the data
%          ((Sigma, omega), Pi_d, Pi_b) and the functor T to such triples
%          (« co-circuits » of a combinatorial 1-complex); page 11 opens a
%          category of topological 1-complexes and breaks off.
%   13-20  Extensions of topological groups 1 -> N -> G -> H -> 1 and base
%          change along u : H' -> H: when pi_0 G' is determined by pi_0,
%          pi_1 data (page 13, the map c onto pi_0 G x_{pi_0 H} pi_0 H' is
%          onto but not always injective); reduction to N discrete, the
%          kernel of c as a quotient of Coker pi_1 u (pages 14-15); for H
%          connected, G is given by phi : pi = pi_1(H) -> z(N) via the
%          universal covering H~; (A) the data (a)-(c) with H~_0, gothic G
%          and strict trivialisations (pages 16-17), Ext(G, N) x
%          Hom_G(pi, z(N)) -> Ext(H mod H~_0, N); page 18, a map of
%          categories (SO, PLO, TUBO, SOB, CO, Kc) with notes on
%          hypercircular complexes and Teichmüller groups T_{g,n}; (B) under
%          alpha) beta) gamma), extensions of H by N correspond to pairs
%          (G~_0, phi) (page 19); explicit construction by contracted
%          product, pi_0(G) = Coker(pi -> G_0), and the start of (B')
%          (page 20), continued on page 21.
%
% Notation fixed for the batch, following batch 2 where they overlap:
% \mathfrak{G} for his curled capital G (a discrete group, pi_0 of H);
% \mathfrak{z}(N) for the curly « 3 » he writes for the centre of N (he says
% so on page 19: « z(N) = centre de N »); \widetilde{H}, \widetilde{H}_0,
% \widetilde{G}_0, \widetilde{H^0} as he writes them; \Pi_1 for fundamental
% groupoids, Cat G for the one-object groupoid; \mathfrak{S}_I for the
% symmetric group; his dotted arrows are `dashed`. Greek letters are always
% in math mode.
%
% Readings to check first: page 2, the parenthesis after pi_0(F)/pi_1(Y)
% (« coinvariants »?) and the slanted margin note; page 3, whether the
% F~ of the lower sequence is Phi; page 6, « Aut(,) »; page 10, the margin
% « catégorie ... topologique des surfaces orientées » and « grossier »;
% page 11, the index of Equ; page 15, the opening sentence (« Montrons que
% ... Coker pi_0 H ... non pas R' ») and « R est un quotient de R' »;
% page 16, the parenthesis after « stricte » and the two cramped
% interlinear lines; page 17, the slanted note in the top-left corner;
% page 18, most of the annotations around the diagram; page 19, the text
% right of the bracket in 1°) and the slanted NB bottom left; page 20, the
% indices of G_! and H~^phi, and « H » in « pi, N, H discrets ».
%
% Counts: \ill 69, \uncertain 83.

\documentclass[11pt,a4paper]{article}
% Le préambule commun à toutes les transcriptions du fonds Grothendieck.
%
% Il tient en une idée : **l'incertitude se compose, elle ne se résout pas.**
% Une transcription qui lisse un mot douteux a détruit la seule chose pour
% laquelle on la faisait — un lecteur doit pouvoir savoir, à chaque endroit,
% ce qui était sur le feuillet et ce qui a été ajouté. Les six macros qui
% suivent sont donc l'appareil critique, et non de la décoration.
%
% Les mêmes commandes sont reconnues par `scripts/render.mjs`, qui produit la
% vue de lecture du site. Une macro ajoutée ici doit l'être aussi là-bas,
% faute de quoi le rendu échouera bruyamment — ce qui est voulu : un rendu
% silencieusement faux est pire qu'un rendu absent.

\ProvidesPackage{grothendieck}[2026/08/08 Transcriptions du fonds Grothendieck]

\usepackage{amsmath,amssymb,amsthm}
\usepackage{mathrsfs}
\usepackage{stmaryrd}
% Les diagrammes commutatifs sont partout dans ce fonds : c'est la langue
% même de la géométrie algébrique de Grothendieck, et une transcription qui
% les rendrait en prose ne transcrirait rien.
\usepackage{tikz-cd}
% Français par défaut — c'est la langue des feuillets — mais l'anglais est
% chargé aussi : la traduction et le résumé sont des documents anglais, et la
% typographie française (espaces avant les deux-points) y serait une faute.
% Ces documents basculent par \selectlanguage{english} en tête de corps.
\usepackage[english,french]{babel}
\usepackage{fontspec}
\usepackage{geometry}
\geometry{a4paper,margin=25mm,marginparwidth=28mm}
\usepackage{marginnote}
\usepackage{xcolor}
% ulem, not soul. `soul`'s \st and \ul are built on a character-by-character
% scan that XeLaTeX's font machinery breaks: the document fails with an
% undefined control sequence rather than degrading. ulem does the same two jobs
% and is Unicode-engine safe. `normalem` stops it hijacking \emph, which the
% transcriptions use for Grothendieck's own emphasis.
\usepackage[normalem]{ulem}
\usepackage{hyperref}
\hypersetup{colorlinks=true,linkcolor=black,urlcolor=black,pdfborder={0 0 0}}

% --- Métadonnées du lot -----------------------------------------------------
% Reprises telles quelles par le script de rendu, qui les affiche en tête de
% la vue de lecture. Elles fixent ce que le fichier prétend couvrir : sans
% elles, une transcription flotte sans point d'attache dans un fonds de seize
% mille feuillets.

\newcommand{\folder}[1]{\gdef\grothFolder{#1}}
\newcommand{\batch}[1]{\gdef\grothBatch{#1}}
\newcommand{\leaves}[2]{\gdef\grothFirst{#1}\gdef\grothLast{#2}}
\newcommand{\foldertitle}[1]{\gdef\grothFoldertitle{#1}}
\gdef\grothFolder{?}\gdef\grothBatch{?}\gdef\grothFirst{?}\gdef\grothLast{?}\gdef\grothFoldertitle{}

% --- Le résumé --------------------------------------------------------------
%
% Il ouvre la lecture modernisée, sous le titre « Résumé », et il est composé
% en petit. Ce n'est pas un ornement : c'est le seul passage du document qui
% ne soit pas des mathématiques, et un lecteur doit pouvoir le reconnaître
% avant de l'avoir lu — puis le sauter s'il connaît déjà le sujet, ou s'y
% arrêter s'il ne le connaît pas. Le filet qui le suit marque la reprise du
% corps.

\newenvironment{resume}
  {\small\setlength{\parskip}{0.45em}}
  {\par\medskip\hrule\medskip}

% --- Mots-clés ---------------------------------------------------------------
%
% En anglais, en clôture du résumé de la lecture modernisée : le vocabulaire
% moderne sous lequel chercher ce que les feuillets construisent. Le manifeste
% du site (`npm run manifest`) les extrait pour étiqueter la cote entière —
% cette ligne est la source unique des tags, il n'y a pas de fichier à côté.
\newcommand{\keywords}[1]{\par\smallskip\noindent
  {\footnotesize\textsc{Keywords} --- \textit{#1}}\par}

% --- L'appareil critique ----------------------------------------------------

% Le numéro de feuillet, en marge. C'est l'ancre qui permet de régler un
% désaccord sur une formule en regardant *un* feuillet plutôt qu'une chemise
% de six cents. Le site s'en sert aussi pour tourner les pages du fac-similé
% quand on fait défiler la transcription.
\newcommand{\leaf}[1]{%
  \marginnote{\footnotesize\color{gray!70}\textsf{#1}}%
  \label{leaf:#1}\ignorespaces}

% Illisible : on ne devine pas. Un mot inventé qui se lit comme les autres est
% le pire résultat possible.
\newcommand{\ill}{\textcolor{red!60!black}{[\,\ldots\,]}}

% Lecture douteuse : le mot est donné, et signalé comme incertain.
\newcommand{\uncertain}[1]{\uline{#1}}

% Ajout de l'éditeur — une lettre manquante, un mot sous-entendu.
\newcommand{\add}[1]{\textcolor{blue!50!black}{[#1]}}

% Biffé par Grothendieck lui-même. Ce qu'il a rayé fait partie du document :
% une rature dit souvent où la pensée a changé de direction.
\newcommand{\struck}[1]{\sout{#1}}

% Note du transcripteur, sur l'état matériel du feuillet.
\newcommand{\note}[1]{\par\smallskip{\small\color{gray!60!black}\textit{[#1]}}\par\smallskip}

% Note marginale de Grothendieck, distincte du corps du texte.
\newcommand{\marginal}[1]{\par\smallskip\noindent\hspace{1em}%
  {\small\color{gray!60!black}#1}\par\smallskip}

% --- Titre ------------------------------------------------------------------

\AtBeginDocument{%
  \begin{center}
    {\large\bfseries Fonds Alexandre Grothendieck --- cote n\textsuperscript{o}~\grothFolder}\\[2pt]
    {\itshape\grothFoldertitle}\\[4pt]
    {\small feuillets \grothFirst--\grothLast\ (lot \grothBatch)}\\[2pt]
    {\footnotesize\color{gray!60!black}Transcription établie d'après le fac-similé numérisé
     par l'Université de Montpellier.\\
     Les lectures incertaines sont soulignées, les illisibles notées [\,\ldots\,],
     les ajouts entre crochets.}
  \end{center}
  \vspace{1em}}

\endinput


\folder{132}
\batch{1}
\pages{1}{20}
\dating{1974}
\watermark{Édition de démonstration}
\foldertitle{Plongements dans surfaces (Ladegaillerie) : notes manuscrites (s.d.), tirés à part (1974).}

\begin{document}

\page{1}
\note{sa pagination : (1)}

\begin{tikzcd}[column sep=small]
X \arrow[dr, "f"'] & & Y \arrow[dl, "g"] \\
& S &
\end{tikzcd}

espaces topologiques, d'où $X \times_S Y$ et
\[
  \Pi_1(X \times_S Y) \xrightarrow{\ \varphi\ } \Pi_1(X) \overset{2}{\times}_{\Pi_1(S)} \Pi_1(Y) .
\]

a) Supposons $f$ une fibration de Serre, alors $\varphi$ induit une bijection
sur les $\pi_0$, un épimorphisme sur les $\pi_1$ associés aux composantes
connexes. I.e.\ $\varphi$ est ess.\ surjectif et surjectif sur les
$\mathrm{Hom}(x, y) \to \mathrm{Hom}(\varphi x, \varphi y)$.

\emph{Démonstration}. Décomposons la situation suivant les composantes
connexes de $X, S, Y$, on est ramené au cas où $S, X, Y$ sont connexes,
non vides, \struck{\ill{}} donc $f$ surjectif. Comme $Y \neq \emptyset$,
prenons $y_0 \in Y$, $x_0 \in X_{s_0}$ où $s_0 = g(y_0)$. Posons
\[
  G = \pi_1(X, x_0), \quad H = \pi_1(Y, y_0), \quad \Sigma = \pi_1(S, s_0) ,
\]
on a alors

\begin{tikzcd}[column sep=small]
G \arrow[dr, "f_1"'] & & H \arrow[dl, "g_1"] \\
& \Sigma &
\end{tikzcd}

permettant de reconstruire mod équivalence
$\Pi_1(X) \overset{?}{\times}_{\Pi_1(S)} \Pi_1(Y)$\note{le signe au-dessus
du $\times$ est ici un « ? », ou un 2 mal formé ; au début de la page et à la
ligne suivante il se lit 2} comme
$\mathrm{Cat}\, G \overset{2}{\times}_{\mathrm{Cat}\, \Sigma} \mathrm{Cat}\, H
=: P$. Les objets de cette dernière sont les objets $\sigma$ de $\Sigma$, et
les hom $\sigma \to \sigma'$ sont les couples $(g, h) \in G \times H$, avec
$g_1(h) \sigma = \sigma' f_1(g)$, i.e.\ $\sigma' = g_1(h) \sigma f_1(g)^{-1}$.
Donc les classes d'isomorphie de $P$ correspondent à l'ensemble des doubles
classes $H \backslash \Sigma / G$, et \struck{le $\pi_1$ correspondant}
$\pi_1(P, \sigma)$ n'est autre que \struck{\ill{}} le stabilisateur de
$\sigma \in \Sigma$ dans $G \times H$, opérant sur $\Sigma$ par
$(h, g) . \sigma = g_1(h) \sigma f_1(g)^{-1}$.
\note{la ligne biffée sous « le stabilisateur » se lit à peu près « aux
composantes des \ill{} » ; « $\pi_1(P, \sigma)$ n'est autre » est ajouté en
interligne au-dessus du passage biffé}

\note{en bas à gauche de la page, un petit carré qui explicite ces flèches}

\begin{tikzcd}
e_\Sigma \arrow[r, "\sigma"] \arrow[d, "f_1(g)"'] & e_\Sigma \arrow[d, "g_1(h)"] \\
e_\Sigma \arrow[r, "\sigma'"'] & e_\Sigma
\end{tikzcd}

\note{au-dessus du carré, « $e_G$ \quad $e_H$ »}

\page{2}
\note{sa pagination : (2)}

Calculons $\pi_0(Z)$ ($Z = X \times_S Y$), en utilisant la fibration de Serre
\[
  F \subset Z \to Y
\]
où $F$ est la fibre de $f : X \to S$, en $s_0$ (avec les points base déjà
choisis, et $z_0 = (x_0, y_0)$) :
\[
  \pi_0(Z) = \pi_0(F) / \pi_1(Y)
\]
(\uncertain{coinvariants} \ill{})
(comme $\pi_0(Y) = 1$) ; d'ailleurs on a, par
la fibration de Serre (et $\pi_0(X) = 1$)\note{la parenthèse « (comme
$\pi_0(Y) = 1$) » est écrite au-dessus de la ligne et rattachée par un
trait}
\[
  F \subset X \to S
\]
\struck{$\pi_1(X) \to \pi_1(S) \to \pi_0(F) \to \pi_0 \ill{} \to 1$}
\struck{i.e.} $\pi_0(F) \simeq \pi_1(S) / \pi_1(X)$ (op.\ à droite), d'où, en composant
\[
  \pi_0(Z) = \pi_1(Y) \backslash \pi_1(S) / \pi_1(X) ,
\]
ce qui montre la bijectivité des $\pi_0$.

Reste. Soit $Z'$ une composante connexe de $Z$. Elle se projette \emph{sur}
$Y$, donc $\exists\, z'_0 \in Z'_{y_0}$. Son image dans $X$ est \struck{bien}
dans $X_{s_0} = F$, donc en la reliant par un arc à $x_0$, on trouve en
appliquant $f$ un lacet en $s_0$, soit $\sigma$ [déterminé mod un élément de
$\mathrm{Im}\, \pi_1(X, x_0) = G$ par chgt.\ de l'arc dans $X$\dots, et
\uncertain{mod.\ de même} un élément de $\mathrm{Im}\, \pi_1(Y, y_0)$ en
\emph{changeant} le point $z'_0$\dots]
\marginal{\uncertain{donne} une description \ill{} $(X, Y)$\dots}\note{note
en biais dans la marge de gauche, à la hauteur des dernières lignes ; « Im »
est ajouté au-dessus de $\pi_1(X, x_0)$}

\page{3}
\note{sa pagination : (3)}

\begin{tikzcd}[column sep=small, row sep=small, nodes={font=\scriptsize}]
\pi_2(Z) \arrow[r] & \pi_2(Y) \arrow[r] \arrow[d] & \pi_1(F) \arrow[r] \arrow[d, "\text{épi}"] & \pi_1(Z) \arrow[r] \arrow[d] & \pi_1(Y) \arrow[r] \arrow[d, "\wr"] & \pi_0(F) \arrow[r] \arrow[d, "\wr"] & \pi_0(Z) \arrow[r] \arrow[d] & 1 \\
& 1 \arrow[r] & \pi_1(\Phi) \arrow[r] & \pi_1(\widetilde{Z}) \arrow[r] & \pi_1(\widetilde{Y}) \arrow[r] & \pi_0(\Phi) \arrow[r] & \pi_0(\widetilde{Z}) \arrow[r] & 1
\end{tikzcd}

\note{sous la seconde ligne, $\pi_1(\widetilde{Y}) \simeq H$ et
$\pi_0(\Phi) \simeq \Sigma / G$ ; $\pi_2(Y)$ est pris entre crochets, et
$\pi_1(Z)$, $\pi_0(Z)$ sont surlignés d'un trait ; le « 1 » qui ouvre la
seconde ligne est repassé}

\begin{tikzcd}[row sep=small]
F \arrow[d, hook] & & \Phi \arrow[d, hook] & \\
X \arrow[d] & & \mathrm{Cat}\, G \arrow[r, "\approx"] \arrow[d] & \widetilde{X} \arrow[d] \\
S & & \mathrm{Cat}\, \Sigma \arrow[r, "\approx"] & \widetilde{S}
\end{tikzcd}

\note{entre les deux colonnes, une colonne biffée : $\pi_1(X) = G$ au-dessus
de $\pi_1(S)$}

$\Phi$ = catégorie des \struck{\uncertain{éléments}} de $\Sigma$, avec
$\mathrm{Hom}_\Phi(\sigma, \sigma')$ = ens.\ des $g \in G$ tels que
$\sigma = \sigma' f_1(g)$ i.e.\ $g \in f_1^{-1}(\sigma'^{-1} \sigma)$ ; son
$\pi_0$ est $\Sigma / G$, ses $\pi_1$ tous isom.\ : $\operatorname{Ker} f_1$.

On a $\Pi_1(F) \to \Phi$, induisant une bijection sur les $\pi_0$, et sur les
$\pi_1$ un épim., dont le noyau est
$\operatorname{Coker}(\pi_2(X) \to \pi_2(S))$

\begin{tikzcd}[column sep=small, row sep=small, nodes={font=\scriptsize}]
\pi_2(X) \arrow[r] & \pi_2(S) \arrow[r] & \pi_1(F) \arrow[r] & \pi_1(X) \arrow[r] \arrow[d, no head, "\wr"] & \pi_1(S) \arrow[r] \arrow[d, no head, "\wr"] & \pi_0(F) \arrow[d, no head, "\wr"] \\
& 1 \arrow[r] & \pi_1(\widetilde{F}) \arrow[r] & \pi_1(\widetilde{X}) \arrow[r] & \pi_1(\widetilde{S}) \arrow[r] & \pi_0(\Phi) \arrow[r] & 1
\end{tikzcd}

\note{la première ligne se poursuivait par « $\to 1$ » et un terme biffé
(\struck{\ill{}} $\pi_0$) ; le $\widetilde{F}$ de la seconde ligne est
peut-être un $\Phi$}

$\sigma = 1$. \quad
$\pi_1(\widetilde{Z}) = \{ (g, h) \mid f_1(g) = g_1(h) \}$

$\pi_1(\widetilde{Z})$ est ext.\ de $g_1^{-1}(\mathrm{Im}\, G)$ par
$\pi_1(\Phi)$ ;
$\pi_1(Z)$ est ext.\ de \uncertain{id.}\ par
$\operatorname{Coker}(\pi_2(Y) \to \pi_1(F))$,
donc est ext.\ de $\pi_1(\widetilde{Z})$ par
$\operatorname{Coker}(\pi_2(X) \times \pi_2(Y) \to \pi_2(S))$

\emph{Donc on trouve} $Z \to \widetilde{Z}$ induit un iso.\ pour les
$\pi_1$ ssi $\pi_2(X) \times \pi_2(Y) \to \pi_2(S)$ est épi.

De plus, $\pi_i(Z) = 0$ pour $i \geq 2$ \struck{$\Leftrightarrow$} dès que
$\pi_i(Y) = \pi_i(F) = 0$ pour $i \geq 2$, or
$\pi_{i+1}(S) \to \pi_i(F) \to \pi_i(X)$ donc il suffit
$\pi_i(X) = \pi_i(Y) = \pi_i(S) = 0$ pour $i \geq 2$, i.e.\ $X, Y, S$ sont
des \ill{}\dots

\page{5}

\begin{tikzcd}
\mathcal{D} \arrow[r, "\varphi"] & \mathcal{C}
\end{tikzcd}

0-objets \quad $(\mathcal{C}, \mathcal{D}, \varphi : \mathcal{D} \to \mathcal{C})$

\note{sous le diagramme : sous $\mathcal{D}$ « groupoïde 1-spécial orienté »,
sous $\mathcal{C}$ « groupoïde 2-spécial », sous la flèche « foncteur
spécial »}

\begin{tikzcd}
\mathcal{D} \arrow[r, "\varphi"] \arrow[d, "u"', "\wr"] & \mathcal{C} \arrow[d, "v"] \\
\mathcal{D}' \arrow[r, "\varphi'"'] & \mathcal{C}'
\end{tikzcd}

1 objets \quad $(u, v, \alpha : \varphi' u \xrightarrow{\sim} v \varphi)$

\note{la 2-flèche $\alpha$ est tracée en diagonale dans le carré ; un
symbole biffé à côté de $v$. Sous $u$ du triplet, une flèche renvoie à
« compatible avec orientations »}

\begin{tikzcd}[column sep=large]
\mathcal{D} \arrow[r, "\varphi"] \arrow[d, bend right=30, "u_1"'] \arrow[d, bend left=30, "u_2"] & \mathcal{C} \arrow[d, bend right=30, "v_1"'] \arrow[d, bend left=30, "v_2"] \\
\mathcal{D}' \arrow[r, "\varphi'"'] & \mathcal{C}'
\end{tikzcd}

2-objets \quad $(\lambda, \mu)$ \quad
$\lambda : u_1 \xrightarrow{\sim} u_2$, $\mu : v_1 \xrightarrow{\sim} v_2$

\note{dans le diagramme, $\lambda$ entre $u_1$ et $u_2$, $\mu$ entre $v_1$ et
$v_2$, et les deux 2-flèches $\alpha_1$, $\alpha_2$ en diagonale,
partiellement repassées}

\begin{tikzcd}
\varphi' u_1 \arrow[r, "\varphi' * \lambda", "\sim"'] \arrow[d, "\alpha_1"', "\wr"] & \varphi' u_2 \arrow[d, "\alpha_2", "\wr"'] \\
v_1 * \varphi \arrow[r, "\mu * \varphi"', "\sim"] & v_2 * \varphi
\end{tikzcd}

commutatif

\note{le premier terme de la ligne du bas est écrit sur une surcharge}

Supposons $\mathcal{C}$ connexe, non vide\add{,} $\mathcal{D}$ non
vide, alors $\varphi : \mathcal{D} \to \mathcal{C}$ est équivalent à un objet
(noté de même) où\note{« non vide » et « $\mathcal{D}$ non vide » ajoutés en
interligne}
\[
  \mathcal{C} = \mathrm{Cat}\, G , \qquad
  \mathcal{D} = \coprod_{i \in I} \mathrm{Cat}\, \mathbb{Z}
\]
donc $\varphi$ donné par des hom $\varphi_i : \mathbb{Z} \to G$, donc par des
$\ell_i = \varphi_i(1) \in G$ \quad ($i \in I$)

idem \quad $\mathcal{C}' = \mathrm{Cat}\, G'$,
$\mathcal{D}' = \coprod_{i' \in I'} \mathrm{Cat}\, \mathbb{Z}$, $\varphi'$
donné par des $\varphi'_{i'} : \mathbb{Z} \to G'$ i.e.\ des
$\ell'_{i'} \in G'$.

$\underline{\mathrm{Hom}}((\mathcal{C}, \mathcal{D}, \varphi);
(\mathcal{C}', \mathcal{D}', \varphi')) \overset{?}{=}$
\struck{\ill{}}

Objets
\begin{itemize}
  \item $\boxed{u : I \xrightarrow{\sim} I'}$ (définit $u$)
  \item $\boxed{v : G \xrightarrow{\sim} G}$ (définit $v$)\note{sic : le
  but est écrit $G$, sans accent}
  \item $\forall i \in I$ \ $\boxed{\alpha(i) \in G'}$ tel que
  \struck{\ill{} $\ell'_{u(i)} = \alpha_i\, v(\ell_i)\, \alpha_i^{-1}$}
  $\boxed{v(\ell_i) = \alpha(i)\, \ell'_{u(i)}\, \alpha(i)^{-1}}$ (ce qui
  définit $\alpha$)
\end{itemize}

\[
  \mathrm{Hom}((u_1, v_1, \alpha_1), (u_2, v_2, \alpha_2)) =
  \left\{ (\lambda, \mu) \;\middle|\;
  \begin{array}{l}
    \lambda \in \mathbb{Z}^I,\ \mu \in G' , \\
    \boxed{v_2 = \operatorname{int}(\mu) \circ v_1} , \\
    \alpha_2(i)\, {\ell'_{u(i)}}^{\lambda_i} = \mu\, \alpha_1(i)
  \end{array}
  \right\}
\]
\marginal{si $u_1 = u_2 =: u$, sinon c'est vide}\note{entouré, et relié par
un trait aux indices de $(u_1, v_1, \alpha_1)$, $(u_2, v_2, \alpha_2)$}
ou aussi
$\boxed{\alpha_2(i)^{-1} \mu\, \alpha_1(i) = {\ell'_{u(i)}}^{\lambda_i}}$
\note{un signe biffé précède $\operatorname{int}(\mu)$ et suit la virgule
après $\mu \in G'$}

\page{6}

Supposons $(\mathcal{C}', \mathcal{D}', \varphi') = (\mathcal{C}, \mathcal{D},
\varphi)$ donc $G = G'$, $I = I'$, $\ell'_i = \ell_i$ ; donc
\uncertain{$\underline{\mathrm{Aut}}(\,,\,)$}
$= \underline{\mathrm{End}}((\mathcal{C}, \mathcal{D}, \varphi))$ est une Gr-catégorie.
\[
  \pi_0(\underline{\mathrm{End}}(\mathcal{C}, \mathcal{D}, \varphi))
  \overset{?}{=}
\]

\textbf{NB} \quad $(u_1, v_1, \alpha_1) \simeq (u_2, v_2, \alpha_2)$
\struck{\ill{}} implique
\begin{enumerate}
  \item[a)] $u_1 = u_2$ $(= u)$
  \item[b)] $v_2$ \emph{conjugué} de $v_1$ i.e.\ $\exists\, \mu \in G'$ tel
  que $v_2 = \operatorname{int}(\mu) v_1$,
\end{enumerate}
d'où
\[
  \begin{array}{rcl}
    v_2(\ell_i) & = & \operatorname{int}(\mu)\, v_1(\ell_i) = \mu\, v_1(\ell_i)\, \mu^{-1} \\
    \| & & \| \\
    \alpha_2(i)\, \ell_{u(i)}\, \alpha_2(i)^{-1} & = & \mu\, \alpha_1(i)\, \ell_{u(i)}\, \alpha_1(i)^{-1} \mu^{-1}
  \end{array}
\]
\note{les deux $\ell$ de la seconde ligne portaient un accent, biffé}
d'où
\[
  (\operatorname{int} \alpha_2(i))\, \ell_{u(i)} =
  (\operatorname{int}(\mu\, \alpha_1(i))) . \ell_{u(i)}
\]
i.e.\ \struck{int} $\alpha_2(i)^{-1} \mu\, \alpha_1(i)$ centralise
$\ell_{u(i)}$.

Or sauf le cas de la « sphère à 1 seul trou », chaque $\ell_j$ est un élément
libre, et \struck{\ill{}} en tous cas\note{« en tous cas » en interligne,
au-dessus du passage biffé} son centralisateur est le groupe engendré par
$\ell_j$ : \struck{dans tous les cas,} donc a) et b) impliquent déjà
l'existence de $\lambda_i$ tel que
$\alpha_2(i)^{-1} \mu\, \alpha_1(i) = \ell_{u(i)}^{\lambda_i}$ (d'ailleurs
unique sauf le cas de la sphère à un seul trou i.e.\ $G = 1$,
$\operatorname{card} I = 1$), donc
\[
  (u_1, v_1, \alpha_1) \simeq (u_2, v_2, \alpha_2) \quad \text{ssi} \quad
  u_1 = u_2 ,\ \overline{v}_2 = \overline{v}_1
\]
(\ill{} au \uncertain{groupe} \ill{}). Donc
$\pi_0(\underline{\mathrm{End}}) \simeq$ groupe des couples
$(u, \overline{v})$, où $u \in \mathfrak{S}_I$,
$\overline{v} \in \operatorname{Autext} G$, tels que $\forall i \in I$, on
ait $\overline{v}(\ell_i)$ conjugué à $\ell_{u(i)}$ (pour $v$ un
\uncertain{aut.}\ de $G$ qui représente $\overline{v}$).

\textbf{NB} \quad $u$ est déjà déterminé par $\overline{v}$, ce qui résulte
du fait que si $i, j \in I$, $i \neq j$, alors $\ell_i$ et $\ell_j$ sont
non conjugués (en fait leurs images dans $G_{\mathrm{ab}}$ sont
distinctes !) Donc $\pi_0(\underline{\mathrm{End}}) \subset
\operatorname{Autext}(G)$, est formé des $\overline{v} \in
\operatorname{Autext}$ tels que $\forall i \in I$, $\overline{v}(\ell_i)$
est conjugué à un des $\ell_j$.
\note{ce NB est marqué d'un trait vertical dans la marge ; la lettre lue
$u$ au début du NB est écrite comme un $\sigma$, et le $\overline{v}$ de la
dernière ligne porte deux barres}

\page{7}

\[
  \pi_1(\underline{\mathrm{End}}) \overset{?}{=}
  \mathrm{Aut}(u = \mathrm{id},\ v = \mathrm{id},\ \alpha = 1 \in G^I) = ?
\]
\note{un signe biffé entre $\alpha$ et $=$}
$\lambda \in \mathbb{Z}^I$, $\mu \in G$ tels que
\begin{enumerate}
  \item[a)] $\mathrm{id}_G = \operatorname{int}(\mu) \circ \mathrm{id}_G$
  i.e.\ $\mu \in Z(G)$ (centre)
  \item[b)] $\forall i$, \ $\mu = \ell_i^{\lambda_i}$
\end{enumerate}

Or, sauf le cas de la sphère à 1 ou 2 trous, le groupe $G$ est libre à
$\geq 2$ générateurs, donc son centre est nul, de plus les $\ell_i$ sont
libres. Donc dans ce cas $\mu = 1$, $\lambda_i = 0$ $\forall i \in I$, donc
$\boxed{\pi_1(\underline{\mathrm{End}}) = 0}$

Dans le cas de la sphère à 2 trous \add{(}i.e.\ le cylindre\add{)}, on a $G \simeq \mathbb{Z}$,
$\ell_i = \pm 1$, \struck{\ill{}} ; les $\ell_i$ sont encore libres (donc les
$\lambda_i$ déterminés par $\mu$) et $G = Z(G)$. Donc
\[
  \pi_1(\underline{\mathrm{End}}) = G \quad (\simeq \mathbb{Z})
\]
(pas canoniquement). Dans le cas de la sphère à 1 trou, $G = 0$, et on trouve
\[
  \pi_1(\underline{\mathrm{End}}) \simeq \mathbb{Z} \quad
  (\text{isom.\ canonique})
\]
\note{« i.e. le cylindre » est ajouté en interligne ; « (pas
canoniquement) » est rattaché par une flèche à « $(\simeq \mathbb{Z})$ »}

\note{un trait horizontal sépare ce qui suit}

Regardons maintenant la 2-catégorie des $(\mathcal{C}, \mathcal{D},
\varphi)$ \emph{avec $\mathcal{D}$ fixe}, donc $u = \mathrm{id}$,
$\lambda = \mathrm{id}$. On trouve pour
$\underline{\mathrm{End}}_{\mathcal{D}\ \mathrm{fixe}}(\mathcal{C}, \mathcal{D},
\varphi)$

objets : couples $(v, \alpha)$ avec
$\left\{ \begin{array}{l}
  v \in \mathrm{Aut}\, G , \\
  \alpha \in G^I \\
  v(\ell_i) = \operatorname{int}(\alpha(i)) . \ell_i
\end{array} \right.$

$\mathrm{Hom}((v_1, \alpha_1), (v_2, \alpha_2))$ = ens.\ des $\mu \in G$ tels
que
$\left\{ \begin{array}{l}
  v_2 = \operatorname{int}(\mu) \circ v_1 \\
  \alpha_2(i) = \mu\, \alpha_1(i)
\end{array} \right.$

On trouve que $(v_1, \alpha_1) \simeq (v_2, \alpha_2)$ \struck{\ill{}}
$\Rightarrow$ \struck{\ill{}} $\overline{v}_1 = \overline{v}_2$\note{« $\overline{v}_1
= \overline{v}_2$ » écrit au-dessus du passage biffé} i.e.
$\exists\, \mu_0$ tel que $v_2 = \operatorname{int}(\mu_0) v_1$, d'où
\[
  \begin{array}{rcl}
    v_2(\ell_i) & = & \underline{\operatorname{int}(\mu_0)\, v_1(\ell_i)} \\
    \| & & \| \\
    \operatorname{int}(\alpha_2(i))\, \ell_i & & \operatorname{int}(\mu_0)\, \operatorname{int} \alpha_1(i)\, v_1(\ell_i)
  \end{array}
\]
\note{sic, $v_1(\ell_i)$ au dernier terme, où l'on attend $\ell_i$}
donc $\alpha_2(i)^{-1} \mu_0\, \alpha_1(i)$ centralise $\ell_i$

\page{8}

\[
  0 \to \mathbb{Z}^I \dashrightarrow TS(\mathcal{C}, \mathcal{D}, \varphi)
  \xrightarrow{\ \Psi\ } T(\mathcal{C}, \mathcal{D}, \varphi) \to
  \mathfrak{S}_I \to 1
\]
\note{sous $TS$ : « Teichmüller spécial » ; sous $T$ : « Teichmüller à
trous » ; sous $\mathfrak{S}_I$ : « groupe symétrique ». Le début de la
suite, avant $\mathbb{Z}^I$, est biffé}
\marginal{\uncertain{exacte} sauf pour la sphère à 1 ou 2 trous}\note{dans la
marge de gauche, sous le début de la suite}

Les $(v, \alpha) \in TS$ dans le noyau de $\Psi : TS \to T$, i.e.\ tels que
$\overline{v} = \mathrm{id}$, sont représentables par des $(v, \alpha)$ avec
$v = \mathrm{id}$, donc
\[
  \ell_i = \operatorname{int}(\alpha(i))\, \ell_i .
\]
Sauf dans le cas de la sphère à 1 trou, on a le centralisateur de $\ell_i$
est le groupe qu'il engendre \emph{et} ce groupe est libre de rang 1, donc
\struck{\ill{}}
\[
  \alpha(i) = \ell_i^{\lambda_i} \qquad \lambda_i \in \mathbb{Z} \text{ bien déterminé.}
\]
On aura
\[
  (\mathrm{id}_G, (\ell_i^{\lambda_i})) \simeq (\mathrm{id}_G, (\ell_i^{\lambda'_i}))
\]
ssi $\exists\, \mu \in Z(G)$ tel que
$\ell_i^{\lambda'_i} = \mu\, \ell_i^{\lambda_i}$.

Dans tous les cas sauf la sphère à 2 trous (cylindre) $Z(G) = 0$, donc la
condition précédente est $(\lambda_i) = (\lambda'_i)$, i.e.\
$\operatorname{Ker} \Psi \simeq \mathbb{Z}^I$. Dans le cas de la sphère à 2
trous, $Z(G) = G \simeq \mathbb{Z}$ \struck{la condition \ill{} que}
$(\ell_1, \ell_2) = (+1, -1)$\note{écrit au-dessus de la ligne biffée}
\struck{$\mu = \ell_i^{\lambda'_i}\, \ell_i^{-\lambda_i}$ \ill{}
$\ell_1^{\lambda_1}, \ell_2^{\lambda_2}$}\note{passage biffé de traits
obliques, lecture partielle}
\ill{} les conditions sur les $\lambda$ (notation additive dans $\mathbb{Z}$)
\[
  \lambda'_1 = \mu + \lambda_1 , \qquad
  -\lambda'_2 = \mu - \lambda_2 \quad \text{i.e.} \quad \lambda'_2 = -\mu + \lambda_2
\]
i.e.
\[
  \operatorname{Ker}(\mathbb{Z}^I \to \operatorname{Ker} \Psi) \simeq
  \operatorname{Ker}(\mathbb{Z}^I \to \mathbb{Z})
\]
\note{sous la dernière flèche, « opération somme »}
donc $\operatorname{Ker} \Psi \simeq \mathbb{Z}$ (canoniquement).

Enfin dans le cas de la sphère à 1 trou, $G = 0$ donc $\alpha = 0$, donc
$TS(\mathcal{C}, \mathcal{D}, \varphi) = 0$.

On trouve d'autre part, trivialement,
\[
  \pi_1(\underline{\mathrm{End}}_{\mathcal{D}\ \mathrm{fixe}}(\mathcal{C},
  \mathcal{D}, \varphi)) = 1 .
\]

\page{10}

\marginal{catégorie \ill{} topologique des surfaces \uncertain{orientées}}\note{écrit
le long du bord gauche, de bas en haut, avec un grand trait vertical qui
embrasse toute la partie supérieure de la page, jusqu'au trait horizontal}

\begin{itemize}
  \item[(a)] $(\Sigma, \omega)$ 1-complexe combinatoire (circulation
  « ordinaire » (ss pts isolés ni bandes isolés)\add{)}, \uncertain{(\ill{})}
  \note{« (\ill{}) », d'un trait plus fin, est ajouté au-dessus et relié par
  un trait à « combinatoire »}
  \item[(c)] $\Pi_d$ $\mathbb{Z}$-groupoïde \struck{\ill{} $\pi_0(\Pi_d)$ fini}
  1-spécial
  \item[(b)] $\Pi_b$ groupoïde 1-spécial (i.e.\ $\pi_0(\Pi_b)$ fini, les
  groupes fond.\ isom.\ à $\mathbb{Z}$)
\end{itemize}
\note{les trois lettres (a), (c), (b) sont entourées, et (a), (c) réunies
par une accolade}

\begin{itemize}
  \item[(a')] $\partial_t(\Sigma, \omega)$ 1-complexe combinatoire
  \struck{\ill{}} \emph{orienté}, dit « les co-circuits » de $\Sigma$.

  $\Pi(\partial_t(\Sigma, \omega))$ \quad $\mathbb{Z}$-groupoïde 1-spécial
  $\xrightarrow{\ \varphi_0\ } \Pi(\Sigma)$

  \item[(b')]
  $\overline{\Pi}_b = \coprod_{i \in B} \Pi_b(i)$ \ (comp.\ connexes),
  \quad $\Lambda(i) = \pi_1 \Pi_b(i)$ (groupe isom.\ à $\mathbb{Z}$)

  $\widetilde{\Pi}_b(i) = \struck{\ill{}} \coprod_{\alpha \in \mathrm{Isom}(\mathbb{Z}, \Lambda(i))}
  \Pi_b(i)$ \quad $\mathbb{Z}$-groupoïde \uncertain{grossier} à deux
  composantes connexes

  \[
    \begin{array}{c}
      \widetilde{\Pi}_b = \coprod_{i \in B} \widetilde{\Pi}_b(i) \\
      \big\downarrow \scriptstyle \varphi_b \\
      \Pi_b
    \end{array}
    \qquad \mathbb{Z}\text{-groupoïde 1-spécial}
  \]
  \note{sous $\overline{\Pi}_b$ et « comp. connexes », un trait souligne
  $\Pi_b(i)$ ; la première barre sur $\Pi_b$ porte un signe peu lisible,
  $\widetilde{\ }$ ou $\sim$ ; « à deux » est écrit au-dessus de
  \struck{\uncertain{connexe}} ; au-dessus de « à deux », un petit signe}

  \item[(c')] $\mathrm{Cat}\, \pi_0(\Pi_d) \xleftarrow{\ \varphi_c\ } \Pi_d$
  \quad $\mathbb{Z}$-groupoïde 1-spécial
\end{itemize}

\[
  T((\Sigma, \omega), \Pi_d, \Pi_b) = (\mathcal{C}, \mathcal{D}, \varphi)
\]
\note{une flèche fine relie $(\mathcal{C}, \mathcal{D}, \varphi)$ à
« $\mathbb{Z}$-groupoïde 1-spécial » de la ligne (c')}
\[
  \begin{aligned}
    \mathcal{D} &= \Pi(\partial_t(\Sigma, \omega)) \amalg \widetilde{\Pi}_b \amalg \Pi_d \\
    \mathcal{C} &= \Pi(\Sigma) \amalg \Pi_b \amalg \mathrm{Cat}(\pi_0(\Pi_d)) \\
    \varphi &= \varphi_0 \amalg \varphi_b \amalg \varphi_d
  \end{aligned}
\]
\textbf{NB} $\varphi$ est foncteur \uncertain{bordant}

\note{un trait horizontal traverse la page}

\begin{itemize}
  \item[(d)] $\mathcal{C}'$ groupoïde
  \item[(e)] $\varphi' : \mathcal{D} \to \mathcal{C}'$
  $(\varphi' = (\varphi'_0, \varphi'_b, \varphi'_d))$ \ foncteur
  \uncertain{bordant}
\end{itemize}

\page{11}

Catégorie \uncertain{non-strictifiée} des 1-complexes topologiques compacts
\ill{} divisés\dots\ combinatoires\note{« combinatoires » est écrit sous la
ligne ; une flèche fine le relie à la ligne suivante}

\begin{itemize}
  \item $K_1$ \quad 1-complexes sans bandes isolées (« sans » défini \ill{}
  \dots)\note{l'indice de $K_1$ est écrit sur une surcharge}
  \item $\Pi_b$ \quad \struck{\ill{}} groupoïde 1-spécial
\end{itemize}
\[
  \underline{\mathrm{Hom}}((K_1, \Pi_b), (K'_1, \Pi'_b)) \simeq
  \mathrm{Cat}(\mathrm{Isom}(K_1, K'_1)) \times
  \underline{\mathrm{Equ}}_{?}(\Pi_b, \Pi'_b)
\]
\note{la formule court jusqu'au bord droit de la feuille, où la fin se
perd ; l'indice de $\underline{\mathrm{Equ}}$, rendu ici par « ? », est illisible ; le reste de la page est blanc}

\page{13}

1) Soit une extension de groupes topologiques
\[
  (E) \qquad 1 \to N \to G \to H \to 1
\]
et un hom.\ de groupes topologiques
\[
  H' \xrightarrow{\ u\ } H
\]
d'où par image inverse une extension $G'$ de $H'$ par $N$ et un hom.\
d'extensions

\begin{tikzcd}[row sep=small]
(E) & 1 \arrow[r] & N \arrow[r] & G \arrow[r] & H \arrow[r] & 1 \\
(E') & 1 \arrow[r] & N \arrow[r] \arrow[u, "\mathrm{id}"] & G' \arrow[r] \arrow[u] & H' \arrow[r] \arrow[u, "u"'] & 1
\end{tikzcd}

\note{le $N$ de la seconde ligne porte une marque en exposant, peut-être un
chiffre}

Passant aux $\pi_0$, \ill{} on en déduit un hom.\ de suites exactes
\struck{(sauf 1 à gauche)}

\begin{tikzcd}[row sep=small, column sep=small]
(S) & \pi_1 H \arrow[r] & \pi_0 N \arrow[r] & \pi_0 G \arrow[r] & \pi_0 H \arrow[r] & 1 \\
(S') & \pi_1 H' \arrow[r] \arrow[u] & \pi_0 N \arrow[r] \arrow[u, no head, "\parallel" description] & \pi_0 G' \arrow[r] \arrow[u] & \pi_0 H' \arrow[r] \arrow[u] & 1
\end{tikzcd}

dont la commutativité du \uncertain{carré} de droite donne lieu à un hom.
\[
  \begin{array}{c}
    \pi_0 G' \xrightarrow{\ c\ } \pi_0 G \times_{\pi_0 H} \pi_0 H' . \\
    \| \\
    \pi_0(G \times_H H')
  \end{array}
\]
\note{la lettre sur la flèche est repassée, lecture $c$ incertaine}

Il est facile de voir sur des exemples que cet hom.\ (qui est toujours
surjectif) n'est pas \uncertain{nécessairement} bijectif. Ainsi, la seule
connaissance de la suite (S) et de $\pi_1(H') \to \pi_1(H)$ ne permet pas, en
général, de \struck{\ill{}} déterminer $\pi_0(G')$.

Notre propos est de préciser des données \uncertain{algébriques} simples,
relatives à l'extension (E), qui permettent de \struck{\ill{}}
déterminer
$\pi_0(G')$ (mod.\ isom.), connaissant $\pi_0 H \to \pi_0 H'$ et
$\pi_1 H \to \pi_1 H'$ induit par $u$.
\note{« déterminer » est en interligne, au-dessus du mot biffé, et
« (mod. isom.) » en interligne, rattaché par un trait à $\pi_0(G')$}
\note{la parenthèse « (discrets) » est ajoutée au-dessus de la ligne, avec un
trait qui la rattache à « la seule connaissance » ; lecture et point
d'attache incertains. Sic : les deux flèches finales vont de $H$ vers $H'$, à
l'inverse de $u$ ; le premier $H$ est écrit sur une surcharge}

\page{14}

$\alpha$) \emph{Réduction au cas $N$ discret.} Observons que $N^0$ est
invariant dans $G$ (et dans $G'$)
\[
  \left\{ \begin{array}{l}
    \pi_0 G = \pi_0 G / N^0 \\
    \pi_0 G' = \pi_0 G' / N^0
  \end{array} \right.
\]
et que \ill{} ; d'ailleurs $G'/N^0$ comme extension de $H$ par
$N/N^0 = \pi_0 N$ \uncertain{est} l'image inverse de $G/N^0$ (extension de
$H$ par $\pi_0 N$) via $H' \xrightarrow{u} H$.
\note{l'accolade et « et que » sont écrits en marge des deux lignes de
formules, sans suite lisible ; sic : « extension de $H$ » pour $G'/N^0$, où
l'on attend $H'$. Sic aussi, dans les deux égalités, $\pi_0 G / N^0$ :
entendre $\pi_0(G/N^0)$}
D'ailleurs le diagramme commutatif

\begin{tikzcd}[row sep=small]
H \arrow[r] & H/N^0 \\
H' \arrow[r] \arrow[u, "u"] & H'/N^0 \arrow[u, "\hat{u}"']
\end{tikzcd}

\note{sic : $H/N^0$, $H'/N^0$ ; on attend $G/N^0$ et $G'/N^0$}
donne un carré commutatif

\begin{tikzcd}[row sep=small]
\pi_i H \arrow[r] & \pi_i(H/N^0) \\
\pi_i H' \arrow[r] \arrow[u, "\pi_i u"] & \pi_i(H'/N^0) \arrow[u, "\pi_i \hat{u}"']
\end{tikzcd}

\note{le $N^0$ du coin inférieur droit est écrit sur une surcharge}

donc la donnée des $\pi_i u$ équivaut à celle des $\pi_i \hat{u}$ (car les
$\to$ horiz.\ sont des isom.) et celle des $\pi_1 u$ \struck{\uncertain{équivaut}}
\uncertain{implique} celle de $\pi_1 \hat{u}$ (car les $\to$ horiz.\ sont des
épi.) \struck{Remplaçons} Divisant tout par $N^0$, on est donc ramené au cas
où $N$ est \emph{discret}. Dans ce cas la suite (S) \uncertain{revient} à
\struck{\ill{}}
\[
  \left\{ \begin{array}{ll}
    \pi_i G \simeq \pi_i H & i \geqslant 2 \\
    \pi_1 G = \operatorname{Ker}(\pi_1 H \to \pi_0 N) & (\pi_0 N = N) \\
    \pi_0 G \text{ est ext.\ de } \pi_0 H \text{ par } \operatorname{Coker}(\pi_1 H \to N)
  \end{array} \right.
\]
\note{au-dessus de la première ligne, un mot biffé et « $i \geqslant 2$ »
écrit en tête ; le but de $\pi_1 H \to$ est $\pi_0 N$ repassé en $N$, avec
$\pi_0 N = N$ en dessous}
et de même pour les $\pi_i G'$. \struck{Notre propos}

Supposons connus \struck{\ill{}} $\pi_i H'$, on connaît \uncertain{donc}
$\pi_i G$\note{sic, pour $\pi_i G'$} si $i \geqslant 2$, et $\pi_1 G'$ comme
$\operatorname{Ker}(\pi_1 H' \to N)$, en notant que \struck{$\pi_1 H' \to$}
\ill{} a comm.\ du

\begin{tikzcd}[row sep=small]
\pi_1 H \arrow[r] & N \\
\pi_1 H' \arrow[r] \arrow[u, "\pi_1(u)"] & N \arrow[u, no head, "\parallel" description]
\end{tikzcd}

donc $\pi_1 H' \to N$ est connu (par $H' \xrightarrow{u} H$) si
$\pi_1 H \to N$ est connu, en termes des seuls $\pi_1 u$. Notre propos est
\struck{donc de} \uncertain{en fait} de déterminer (pour $H'$ et
$u : H' \to H$) l'extension
\note{la page s'arrête sur ce mot ; la suite est en page 15}

\page{15}
\note{suite de la page 14}

$\pi_0 G'$ de $\pi_0 H'$ par $R' = \operatorname{Coker}(\pi_1 H' \to N)$.
\ill{} \uncertain{Montrons} que $\pi_0 G$ est lui-même une extension de
\uncertain{$\operatorname{Coker}$} $\pi_0 H$ par
$\operatorname{Coker}(\pi_1 H \to N)$, son image inverse par
$\pi_0 H' \to \pi_0 H$ est une extension de $\pi_0 H'$ par
$R = \operatorname{Coker}(\pi_1 H \to N)$ \struck{\ill{}}, \uncertain{non pas}
$R'$.
\note{« Coker », souligné, porte en dessous
$\pi_0 G \times_{\pi_0 H} \pi_0 H'$, relié par un trait à « image
inverse » : c'est ce produit fibré qui est l'image inverse. Le mot
« Coker » devant $\pi_0 H$ est de lecture incertaine}
\uncertain{$R$ est} un quotient de $R'$ par
\[
  \frac{\operatorname{Im}(\pi_1 H \xrightarrow{\varphi} N)}
       {\operatorname{Im}(\pi_1 H' \xrightarrow{\varphi'} N)}
  \simeq \pi_1 H / (\operatorname{Ker} \varphi + \operatorname{Im} \pi_1(u))
\]
donc on a
\[
  1 \to \pi_1 H / (\operatorname{Ker} \varphi + \operatorname{Im} \pi_1(u))
  \to \pi_0 G' \xrightarrow{\ c\ } \pi_0 G \times_{\pi_0 H} \pi_0 H' \to 1 .
\]
\note{le $R$ initial porte un petit signe en exposant ; le terme
$\pi_1 H / (\operatorname{Ker} \varphi + \operatorname{Im} \pi_1(u))$ est
écrit sous la ligne précédente, à gauche, et renvoyé par un trait au membre
de droite de l'isomorphisme}

Cel\add{a} \struck{\uncertain{montre}} \uncertain{met} en évidence de façon
précise comment le noyau de $c$ est lié au $\pi_1 H$ et $\pi_1 u$, c'est un
quotient de $\operatorname{Coker}\, \pi_1 u$.

$\beta$) \emph{Cas $H$ connexe.} Alors l'ext.\ $G$ est \uncertain{canonique},
à isom.\ unique près, par un hom.
\[
  \struck{\pi \to}\ \pi \overset{\mathrm{déf}}{=} \pi_1(H)
  \xrightarrow{\ \varphi\ } \mathfrak{z}(N)
\]
et peut se décrire \ill{} à l'aide du groupe \struck{\ill{}} universel
$\widetilde{H}$ de $H$ (ext.\ de $H$ par $\pi$) en prenant son « image
directe » par $\pi \xrightarrow{\varphi} N$. Alors l'extension $G'$ de $H'$
par $N$ peut s'obtenir en prenant l'image inverse
$H' \times_H \widetilde{H}$ de l'ext.\ $\widetilde{H}$ de $H$ par $\pi$ via
$H' \to H$, et prenant l'image par
$\pi \xrightarrow{\varphi} \mathfrak{z}(N) \hookrightarrow N$. Il n'y a
\uncertain{pas} de \uncertain{problème} si $H'$ \struck{\ill{}} n'est pas
connexe (\uncertain{décrivons} les ext.\ de $\pi_0 H'$ \uncertain{qui}
sont \uncertain{\ill{}} !) --- \uncertain{p.\ ex.} $N$ discret.
\struck{discret (plus généralement \ill{})}
\struck{Il faut prendre l'\ill{}}
\note{deux lignes biffées closent la page ; le « $\mathfrak{z}$ » rend le
« 3 » bouclé qu'il écrit pour le centre, comme en page 21. Au-dessus de
$\pi \xrightarrow{\varphi} \mathfrak{z}(N)$ le signe
$\overset{\mathrm{déf}}{=}$ porte « déf » ; « par » dans « l'image inverse
\dots\ de l'ext. » est écrit sur une surcharge}

\page{16}

\begin{itemize}
  \item[(a)] $1 \to N \to G \to H \to 1$ \quad ext.\ de groupes top.\
  [$N$ discret]
  \item[(b)] $1 \to \pi \to \widetilde{H} \to H \to 1$ \quad ext.\ de
  groupes top.\ [$\pi$ discret]
  \item[(c)] $1 \to \widetilde{H}_0 \to \widetilde{H} \to \mathfrak{G} \to 1$
  \quad ext.\ de gr.\ top.\ [$\mathfrak{G}$ discret]
\end{itemize}
\note{au-dessus de (a), un « (A) » entouré ; devant (b), un signe biffé}

On suppose données (b) et (c) et $N$, on s'intéresse aux ext.\ $G$ de $H$ par $N$, \emph{munies}
d'une trivialisation \emph{stricte} (\uncertain{i.e.\ comm.\ avec prod.\
produits}) de l'image inverse de cette extension par
$\widetilde{H}_0 \to H$ (composé $\widetilde{H}_0 \hookrightarrow
\widetilde{H} \to H$) i.e.\ d'un hom.\ $\widetilde{H}_0 \to G$ qui relève
$\widetilde{H}_0 \to H$ et centralise $N$
\note{« et $N$ » est en interligne, relié par un trait ; la parenthèse
après « stricte » est ajoutée au-dessus de la ligne, sa lecture est très
incertaine}
(\textbf{NB} $N$ étant discret, si $\widetilde{H}_0$ est connexe et
simplement connexe, il existe pour $G$ donné une unique telle
trivialisation\dots)
\uncertain{ce qui} \ill{} \uncertain{de plus} (par ce splitting on
\uncertain{associe} \ill{} \ill{} de $\widetilde{G}$ sur
$\widetilde{H}_0$ (\ill{}) et sur $\mathfrak{G}$ (via
$\widetilde{G} \to G$))
\note{deux lignes serrées, ajoutées entre les lignes, dont l'essentiel ne se
lit pas}

Soit $\widetilde{G}$ l'image inverse de $G$ par $\widetilde{H} \to H$
\marginal{(cf ci-dessous bien déf.\ de $\widetilde{G}$)}
\[
  1 \to N \to \widetilde{G} \to \widetilde{H} \to 1
\]
la donnée d'une ext.\ $G$ est équivalente à celle de $\widetilde{G}$,
\uncertain{munie} d'une trivialisation \emph{stricte} sur $\pi$. La \struck{trivialisation de}
\uncertain{structure} supplémentaire sur l'ext.\ $G$ (trivialisation de
l'image inverse par $\widetilde{H}_0 \to H$) équivaut donc à une
trivialisation \emph{stricte} de $\widetilde{G}$ sur le
sous-groupe $\widetilde{H}_0$. Dans ces conditions on doit étudier les
ext.\ $\widetilde{E}$ de $\widetilde{H}$ par $N$, trivialisées
\emph{strictement} \uncertain{simultanément} sur $\pi$ et sur
$\widetilde{H}_0$ [rappelons que cette dernière structure, \struck{\ill{}}
sur $\widetilde{H}_0$, est automatique dans les conditions dites]. Or une
ext.\ \struck{\ill{}} $\widetilde{E}$ de $\widetilde{H}$ par $N$
trivialisée \emph{strict}\add{ement} sur $\widetilde{H}_0$
équivaut à une ext.\ de $\mathfrak{G}$ par $N$, soit $\widetilde{G}_0$ :
\[
  1 \to N \to \widetilde{G}_0 \to \mathfrak{G} \to 1
\]
\note{au-dessus de $\widetilde{G}_0$, un signe biffé}
($\widetilde{E}$ est l'image inverse par $\widetilde{H} \to \mathfrak{G}$).
De plus, il faut donner un splitting \emph{strict} des diagrammes
\note{à droite de « Soit $\widetilde{G}$ \dots\ par $\widetilde{H} \to H$ »,
la parenthèse « (cf ci-dessous \dots) » est reliée par un trait ; chacun des
mots « stricte », « strictement », « strict » en italique est ajouté en
interligne, relié par un trait}
\note{la page s'arrête sur ce mot ; la suite, pages 17 à 20, puis 21}

\page{17}
\note{suite de la page 16}

\begin{tikzcd}[row sep=small]
1 \arrow[r] & N \arrow[r] & \widetilde{G}_0 \arrow[r] & \mathfrak{G} \arrow[r] & 1 \\
& & & \pi \arrow[u] \arrow[ul, dashed, "\varphi"] &
\end{tikzcd}

\note{la flèche $\varphi$ est pointillée ; au-dessus de $\widetilde{G}_0$,
un signe biffé}

\uncertain{i.e.\ un} \uncertain{relèvement}, et \uncertain{compatible}
($\widetilde{G} \to \widetilde{G}_0$), pour $\widetilde{G}$ \ill{} sur
$\pi$ (via $\widetilde{H}$ et $\widetilde{G} \to \widetilde{H}$) et
\struck{\ill{}} $\widetilde{G}_0$ (via $G$\dots
\note{le début de la ligne est surchargé ; « et
$\widetilde{G} \to \widetilde{G}_0$ » est écrit sous la ligne, et
$\widetilde{G}_0$, à la fin, au-dessus d'un mot biffé}

\emph{Cas particulier} : $\pi \subset \widetilde{H}_0$, i.e.\
$\pi \to \mathfrak{G}$ est trivial, alors la donnée de $\varphi$ équivaut à
la donnée d'un hom.
\[
  \varphi : \pi \to \mathfrak{z}(N) .
\]
\marginal{si $\pi$ \uncertain{connexe} et discret, l'équivalence de $G$ par
$\pi$ se fait via $\pi_0(H)$ (ou $\pi_0(\widetilde{H})$) et celle de
$\mathfrak{G}$ par $\mathfrak{z}(N)$ se fait via \struck{\ill{}}
$\widetilde{H}$ \uncertain{trivial} et si $\mathfrak{G}$ discret et
$\widetilde{H}_0$ connexe, \uncertain{$\pi_0(H)$}
$= \pi_0(\widetilde{H}) = \mathfrak{G}$ \ill{} dans le \uncertain{cas} de
$\varphi : \pi \to \mathfrak{z}(N)$ \uncertain{et} $\mathfrak{G}$-hom.}
\note{note écrite en biais dans l'angle supérieur gauche, à côté des lignes
qui précèdent, et reliée à « Cas particulier » par un trait ; lecture très
partielle}

\note{un trait horizontal traverse la page. En dessous, jusqu'au second
trait, une série de diagrammes de brouillon, barrés de longs traits
obliques : les suites $1 \to \pi \to \widetilde{H} \to H \to 1$ et
$1 \to \pi \to \widetilde{H}' \to H' \to 1$ reliées par
$f : \widetilde{H}' \to \widetilde{H}$, avec $\widetilde{H}'_0 =
f^{-1}(\widetilde{H}_0)$ ; un carré de suites
$1 \to \pi \to \widetilde{H} \to H \to 1$ au-dessus de
$\to \widetilde{H}_0 \to G \to 1$ et d'une ligne $N \to N$, avec une flèche
$\varphi$ en biais ; une grille de sous-groupes $H_1 \cap H_2 \subset H_1,
H_2 \subset H$ et de leurs quotients ; enfin les suites
$1 \to \pi \to \widetilde{H} \to H \to 1$ et
$1 \to \pi \to \widetilde{H}_0 \to H_0 \to 1$ sous
$\mathfrak{G} = \mathfrak{G}$. Ces brouillons ne sont pas repris dans la
suite}

\[
  \operatorname{Ext}(\mathfrak{G}, N) \times
  \operatorname{Hom}_{\mathfrak{G}}(\pi, \mathfrak{z}(N))
  \xrightarrow{\ \sim\ } \operatorname{Ext}(H \bmod \widetilde{H}_0, N)
  \to \operatorname{Ext}(H, N)
\]
\note{une flèche verticale descend de $\operatorname{Ext}(\mathfrak{G}, N)$
vers « équivaut à » ; $\operatorname{Ext}(H, N)$ est écrit sous
$\operatorname{Ext}(H \bmod \widetilde{H}_0, N)$}

\[
  1 \to N \to E \to \widetilde{H} \to 1
  \quad \text{équivaut à} \quad H^2(\mathfrak{G}, N)
\]
\note{sous $E$ et $\widetilde{H}$, des $H$ à demi effacés ; à gauche un
mot biffé}

espace homogène sous $H^2(H, \mathfrak{z}(N))$

\[
  H^2(\mathfrak{G}, \mathfrak{z} N) \to H^2(H, \mathfrak{z} N)
\]
\uncertain{via} $H \to \mathfrak{G}$
\[
  \operatorname{Ext}(\mathfrak{G}, N) \times \operatorname{Hom}(\pi, N)
  \to H^2(H, \mathfrak{z}(N)) .
\]
\note{la diagonale qui barre les brouillons traverse aussi ces dernières
lignes, sans qu'il soit sûr qu'elle les annule}

\page{18}

\note{page organisée autour d'un schéma de sigles reliés par des flèches ;
il est redessiné ici, les annotations qui l'entourent sont données
ensuite}

\begin{tikzcd}[column sep=small, row sep=small]
& \mathrm{SO} & & & \\
& & \mathrm{PLO} \arrow[ul] \arrow[dl] \arrow[drr] & & \\
& \mathrm{TUBO} \arrow[dl] \arrow[r] & \mathrm{SOB} \arrow[dr, "\partial"] & & \mathrm{SOB} \arrow[dl, "-\partial\ (\simeq \partial)"] \\
\mathrm{Kc} & & & \mathrm{CO} &
\end{tikzcd}

\note{sous $\mathrm{PLO}$ : « (avec $\pi_0(K) \to \pi_0(S)$ surjectif) » ;
à côté du second $\mathrm{SOB}$ : « \uncertain{y} \uncertain{circuits} » ;
sous $\mathrm{CO}$ : « $\simeq$ \uncertain{Equiv} $(\mathbb{N},
\{\mathfrak{G}_n\})$ » ; la flèche $\mathrm{TUBO} \to \mathrm{Kc}$ porte un
signe surchargé, et une flèche courbe va de $\mathrm{PLO}$ vers
« \uncertain{classif.} » à gauche}

Accolade à droite de $\mathrm{PLO}$ :

complexes hypercirculaires (i.e.\ \uncertain{complexes} \struck{\ill{}}
circulaires $K$ avec une partition $\Pi'$ de $\mathrm{Circ}(K)$ (ens.\ des
circuits) et applications $\Pi' \xrightarrow{g} \mathbb{N}$.

extensions des \emph{groupes} de symétries de 1-complexes combinatoires
hypercirculaires $\Pi'$ \struck{\ill{}} des groupes de Teichmüller
\uncertain{spéciaux} $\Pi_{g_i, n_i}$ ; ces ext.\ sont connues
\uncertain{quand} on connaît les $\Pi'_{g_i, n_i}$ \uncertain{séparément}, et
les hom.\ canoniques
\[
  \mathbb{Z}^{n_i} \xrightarrow{\ \lambda\ } \operatorname{Centre}(L_{g_i, n_i})
\]
et les ext.\ de Teichmüller
\[
  1 \to T_{g_i, n_i} \to L_{g_i, n_i} \to \mathfrak{G}_{n_i} \to 1
\]
\note{au-dessous, un passage encadré, entièrement hachuré, où l'on devine
$L_{g_i, n_i} \to T_{g_i, n_i}$ et « de noyau $\operatorname{Im}(\mathbb{Z}^{n_i}
\to L_{g_i, n_i})$ » ; il est remplacé par « \dots, \ill{} ce qui est dit
dessus », avec « en \uncertain{donnant} \ill{} » en interligne. Une flèche
courbe y renvoie depuis l'accolade}

Sous $\mathrm{SOB}$ (à gauche) :
familles finies de couples $(g_i, n_i)$ (genre, nb.\ de trous)
$g \geqslant 0$, $n \geqslant 0$ ; « groupes de Teichmüller à trous » :
familles finies de groupes extérieurs à lacets. Tout se détermine : à l'aide
des « gr.\ de Teichmüller à trous » ordinaires $T_{g,n}$.

Sous $\mathrm{Kc}$ :
$\simeq$ [1-complexes combinatoires circulaires] ;
\uncertain{groupoïdes} des classes d'isomorphie de 1-complexes circulaires
connexes, groupes de symétries des complexes circulaires : tout se ramène à
la détermination des groupes de symétries des 1-complexes circulaires
connexes.

\note{une flèche courbe part de $\mathrm{TUBO}$ et mène à ce qui suit}
$\to$ il suffit de le déterminer pour les \struck{\ill{}} cas connexes
\uncertain{i.e.} les 1-compl.\ \uncertain{circ.}\ connexes. Le calcul des
$g_i n_i$ est donné par Yves. Il reste à \uncertain{voir} \ill{} les
$K_g$ circulaires\note{« circulaires » en interligne, au-dessus d'un mot
biffé} de déterminer
\[
  \boxed{\operatorname{Aut}_{\mathrm{Circ}}(K) \to \operatorname{Aut}(T_{?}(K)) \simeq T_{g,n}}
\]
\note{l'indice rendu par « ? » est illisible (\ill{}) ; sous la première
flèche, un signe biffé. Le prénom « Yves » se lit nettement}

\page{19}

\marginal{(B)}\note{entouré, en haut à gauche}

Supposons à la fois
\begin{itemize}
  \item[$\alpha$)] $\widetilde{H}_0$ est connexe et simplement connexe
  \item[$\beta$)] $\pi \subset \widetilde{H}_0$
  \item[$\gamma$)] $\pi$, $N$, $\mathfrak{G}$ discrets
\end{itemize}

Comme $\mathfrak{G}$ discret, $\widetilde{H}_0$ connexe, on a
$\left\{ \begin{array}{l}
  \widetilde{H}_0 = \widetilde{H}^0 \\
  \mathfrak{G} = \pi_0 \widetilde{H}
\end{array} \right.$ (\uncertain{composante} \uncertain{neutre}).
Comme $\pi \subset \widetilde{H}_0$, on a que $H = \widetilde{H}/\pi$ est une
extension de $\mathfrak{G}$ par $H^0 = \widetilde{H}_0/\pi$
\[
  \boxed{H^0 = \widetilde{H}_0/\pi , \quad \mathfrak{G} = \pi_0(H)}
\]
Comme $\pi$ discret\add{,} $\widetilde{H}_0$ est simplement connexe,
\struck{c'est le} le revêtement \emph{universel} de
$H^0 = \widetilde{H}_0/\pi$
\note{« $\pi$ discret, » et « $\widetilde{H}_0$ est » en interligne}
\[
  \boxed{\left\{ \begin{array}{l}
    \widetilde{H^0} = \text{rev.\ universel de } H^0 \\
    \pi = \pi_1(H^0) = \pi_1(H)
  \end{array} \right.}
\]

Donc la situation équivaut à :

1°) donnée d'un groupe top.\ $H$, \struck{\ill{}} i.e.\ [ce qui est
\uncertain{donc} ext.\ d'un \emph{groupe} discret $\mathfrak{G} = \pi_0(H)$
par \struck{\ill{}} un groupe connexe $H^0$ ; en donnant \struck{\ill{}} à un
groupe \uncertain{\ill{}} universel $\widetilde{H^0}$, extension de $H^0$
par $\pi \simeq \pi_1(H^0)$]
\note{à droite du crochet, « \uncertain{groupe} \struck{connexe} à un
\uncertain{groupe} $\mathfrak{G}$ \uncertain{opérant} », sans suite sûre}

2°) Une extension $\widetilde{H}$ de $H$ par $\pi$ \uncertain{qui}
prolonge l'extension $\widetilde{H^0}$ de $H^0$ par $\pi = \pi_1(H^0)$.

\textbf{NB} Une telle extension n'est pas nécessairement unique
\uncertain{à} \uncertain{isom.}\ près, elle est unique
\struck{\ill{}} (i.e.\ mod.\ action de $H^2(\mathfrak{G}, \pi)$)\note{en
interligne, au-dessus du passage biffé}
comme extension de $\mathfrak{G}$ par $\pi$, et sauf erreur l'obstruction à
l'existence d'une telle ext.\ est dans $H^3(\mathfrak{G}, \pi)$. Une telle
extension existe si $H$ est produit semi-direct $\mathfrak{G} . H^0$ ; car
on peut prendre alors $\widetilde{H} = \mathfrak{G} . \widetilde{H^0}$.

3°) Le groupe discret $N$.

Avec ces données, on trouve que \emph{la catégorie des extensions de $H$
par $N$} équivaut à celle des couples $(\widetilde{G}_0, \varphi)$, où
$\widetilde{G}_0$ est une ext.\ de $\mathfrak{G}$ par $N$ et
$\varphi : \pi \to \mathfrak{z}(N)$ un $\mathfrak{G}$-hom.\
($\mathfrak{z}(N)$ = centre de $N$).
\note{au-dessus de $\widetilde{G}_0$, un signe biffé, et un $G$ au-dessus de
« \dots de $H$ » ; le « $\mathfrak{G}$- » de « $\mathfrak{G}$-hom. » est
écrit sur une surcharge}

\marginal{\textbf{NB} $G_0 = \pi_0 \widetilde{G} = \widetilde{G}/\widetilde{G}^0$
(\struck{\ill{}} $\widetilde{G}^0 \to \widetilde{H}^0$ \uncertain{est}
\uncertain{un} \uncertain{isom.}) comme $G = \widetilde{G}/\pi$, et
\uncertain{on} \uncertain{a} \dots\ $\to \pi_0 G \to \pi_0 G \to 1$}
\note{écrit en biais dans l'angle inférieur gauche, avec le diagramme
$\pi \to \pi_0 G = G_0 \supset N \supset \mathfrak{z}(N)$, une flèche
directe $\pi \to \mathfrak{z}(N)$ ; lecture partielle}

\page{20}

On peut aussi expliciter ainsi la construction de l'extension $G$
\add{(}de $H$ par $N$\add{)} associée à l'ext.\ $\widetilde{G}_0$ de $\mathfrak{G}$ par $N$ et le
$\mathfrak{G}$-hom.\ $\varphi : \pi \to \mathfrak{z}(N)$ :
\note{« de $H$ par $N$ » en interligne ; devant $\pi$, un signe biffé ;
au-dessus de $\widetilde{G}_0$, un autre}

Soit $G_!$ l'ext.\ de $H$ par $N$ image inverse de l'ext.\ $G_0$ (de
$\mathfrak{G}$ par $N$) via $H \to \mathfrak{G}$. Considérons l'extension
$\widetilde{H}$ de $H$ par $\pi$, le $\mathfrak{G}$-hom.\
$\varphi : \pi \to \mathfrak{z}(N)$ nous donne une ext.\ \uncertain{centrale}
$\widetilde{H}^{\varphi}_{\ast}$ de $H$ par $\mathfrak{z}(N)$, correspondant
aux opérations de $H$ (via $\mathfrak{G}$) sur $\mathfrak{z}(N)$ par l'ext.\
$G_!$. On peut donc prendre le produit contracté des deux
extensions $G_!$ et
$\widetilde{H}^{\varphi}$, on trouve le $G$ cherché.
\note{« extensions » en interligne ; l'indice de $G_!$ est un point d'exclamation, avec au-dessus un
petit signe qui peut être un $\varphi$ ; l'indice de
$\widetilde{H}^{\varphi}_{\ast}$ est incertain}

D'ailleurs, \struck{\ill{}} comme $G$ est un revêtement de $H$, on aura
\[
  \begin{array}{ll}
    \pi_i(G) \xrightarrow{\ \sim\ } \pi_i(H) & \text{si } i \geqslant 2 \\
    \pi_1(G) = \operatorname{Ker}(\pi_1(H) \to \pi_0(N)) &
  \end{array}
\]

\begin{tikzcd}[row sep=small]
\pi_1(H) \arrow[r] \arrow[d, no head, "\parallel" description] & \pi_0(N) \arrow[d, no head, "\parallel" description] \\
\pi \arrow[r] \arrow[dr, "\varphi"'] & N \\
& \mathfrak{z}(N) \arrow[u, hook]
\end{tikzcd}

\begin{tikzcd}[column sep=small, row sep=small]
\pi_1(H) \arrow[r] \arrow[d, no head, "\parallel" description] & N \arrow[r] & \pi_0(G) \arrow[r] & \pi_0(H) \arrow[r] \arrow[d, no head, "\parallel" description] & 1 \\
\pi \arrow[r, "\varphi"'] & \mathfrak{z}(N) \arrow[u, hook] & & \mathfrak{G} &
\end{tikzcd}

i.e.\ $\pi_0 G$ est une ext.\ de $\mathfrak{G}$ par
$\operatorname{Coker}(\varphi : \pi \to N)$, et on vérifie que l'on a une
suite exacte
\[
  \pi \to G_0 \to \pi_0(G) \to 1
\]
\note{au-dessous, une ligne entièrement hachurée où se lit
$\pi_0(G) \to \pi_0(\dots)$}
(i.e.\ $\pi_0(G) \simeq \operatorname{Coker}(\pi \to G_0)$) où
$\pi \to G_0$ est le composé
\[
  \pi \to \mathfrak{z}(N) \hookrightarrow N \hookrightarrow G_0 \qquad \dots
\]

\marginal{(B')}\note{entouré, dans la marge gauche}
Supposons seulement $\pi, N, H$ discrets, $\pi \subset \widetilde{H}_0$
(conditions $\beta$) et $\gamma$)) mais pas nécess.\ $\alpha$)) (i.e.\ que
$\widetilde{H}_0$ soit connexe et simplement connexe).
\note{le $H$ de « $\pi, N, H$ discrets » est un $H$ capital appuyé ; la
condition $\gamma$) de la page 19 porte sur $\mathfrak{G}$, qu'on attend
peut-être ici}
$N$ étant discret, on a encore
\[
  \pi_i(G) \xrightarrow{\ \sim\ } \pi_i(H) \quad \text{si } i \geqslant 2
\]
\marginal{TSVP}\note{en bas à droite ; la suite est en page 21}

\end{document}
